\section*{Hoofdstuk \ref{ch:g7:priorities} --- Voorrangsregels}

\begin{solution}{exo:g7:priorities:1}
$5 + 3 \times 6 = 5 + 18 = 23$.

$(5 + 3) \times 6 = 8 \times 6 = 48$.

$30 - 12 \div 4 = 30 - 3 = 27$.

$(30 - 12) \div 4 = 18 \div 4 = 4.5$.
\end{solution}

\begin{solution}{exo:g7:priorities:2}
$24 - 6 + 2 = 18 + 2 = 20$ (van links naar rechts).

$24 \div 6 \times 2 = 4 \times 2 = 8$ (van links naar rechts).

$24 - 6 \times 2 = 24 - 12 = 12$ (eerst de vermenigvuldiging).

$24 \div (6 \times 2) = 24 \div 12 = 2$.
\end{solution}

\begin{solution}{exo:g7:priorities:3}
$7 + 3 \times (10 - 6) = 7 + 3 \times 4 = 7 + 12 = 19$.

$40 - (2 + 3) \times 4 = 40 - 5 \times 4 = 40 - 20 = 20$.

$60 \div \bigl(2 \times (7 - 4)\bigr) = 60 \div (2 \times 3) = 60 \div 6 = 10$.
\end{solution}

\begin{solution}{exo:g7:priorities:4}
$\dfrac{18 + 6}{8} = \dfrac{24}{8} = 3$; \quad
$\dfrac{30}{7 - 2} = \dfrac{30}{5} = 6$; \quad
$\dfrac{5 \times 8}{4 + 6} = \dfrac{40}{10} = 4$.
\end{solution}

\begin{solution}{exo:g7:priorities:5}
$8 \times 102 = 8 \times (100 + 2) = 800 + 16 = 816$.

$25 \times 99 = 25 \times (100 - 1) = 2500 - 25 = 2475$.

$14 \times 12 + 14 \times 8 = 14 \times (12 + 8) = 14 \times 20 = 280$.

$37 \times 45 - 37 \times 35 = 37 \times (45 - 35) = 37 \times 10 = 370$.
\end{solution}

\begin{solution}{exo:g7:priorities:6}
``Het \omterm{ex:g1:subtraction:difference}{verschil} van $50$ en het \omterm{def:g2:mult:def}{product} van $6$ en $7$'':
$50 - 6 \times 7 = 50 - 42 = 8$ (geen haakjes nodig; de voorrangsregel doet het
werk).

``Het \omterm{def:g2:mult:def}{product} van de \omterm{def:g1:addition:def}{som} van $4$ en $5$ met het \omterm{ex:g1:subtraction:difference}{verschil} van $10$ en $8$'':
$(4 + 5) \times (10 - 8) = 9 \times 2 = 18$.
\end{solution}

\begin{solution}{exo:g7:priorities:7}
$3 \times 2.50 + 2 \times 1.20 = 7.50 + 2.40 = 9.90$ euro.
\end{solution}

\begin{solution}{exo:g7:priorities:8}
(a) zonder haakjes: $4 + 2 \times 5 - 1 = 4 + 10 - 1 = 13$.

(b) $(4 + 2) \times 5 - 1 = 30 - 1 = 29$.

(c) $(4 + 2) \times (5 - 1) = 6 \times 4 = 24$.
\end{solution}

\begin{solution}{exo:g7:priorities:9}
$k = 1$: $12 + 4 \times 1 = 16$ en $16 \times 1 = 16$ --- die komen overeen (en
daarom is $k = 1$ een slechte test!). $k = 2$: $12 + 4 \times 2 = 20$ maar
$16 \times 2 = 32$: verschillend. De twee uitdrukkingen zijn dus in het algemeen
\emph{niet} gelijk. De fout $12 + 4 \times k = 16 \times k$ komt van eerst
$12 + 4$ optellen, en zo de voorrang van $\times$ negeren.
\end{solution}

\begin{solution}{exo:g7:priorities:10}
\emph{1.} Groepsprijs: $12 + 6 \times n$. Voor $n = 5$: $12 + 30 = 42$ euro.

\emph{2.} Losse prijs: $9n$. Vergelijk: $n = 3$: $30$ tegen $27$ (los
goedkoper); $n = 4$: $36$ tegen $36$ (gelijk); $n = 5$: $42$ tegen $45$ (groep
goedkoper). Vanaf $5$ personen wint het groepstarief (elke extra persoon kost
$6$ in plaats van $9$, dus wordt het \omterm{ex:g1:subtraction:difference}{verschil} alleen groter).
\end{solution}

\begin{solution}{exo:g7:priorities:11}
Mogelijke antwoorden (er zijn er veel):
\[
24 = 1 \times 2 \times 3 \times 4, \qquad
10 = 1 + 2 + 3 + 4, \qquad
1 = (4 - 3) \times (2 - 1).
\]
\end{solution}

\begin{solution}{pb:g7:priorities:1}
\textbf{1.} $7 \times 102 = 7 \times (100 + 2) = 700 + 14 = 714$; \quad
$9 \times 98 = 9 \times (100 - 2) = 900 - 18 = 882$.

\textbf{2.} $17 \times 13 + 17 \times 7 = 17 \times (13 + 7) = 17 \times 20 =
340$; \quad $45 \times 101 - 45 = 45 \times (101 - 1) = 45 \times 100 =
4\,500$.

\textbf{3.} $5 \times 87 = 870 \div 2 = 435$. Waarom het mag: $5 = 10 \div 2$,
dus is met $5$ vermenigvuldigen hetzelfde als met $10$ vermenigvuldigen en dan
door $2$ delen (de volgorde van $\times$ en $\div$ op hetzelfde niveau, van
links naar rechts, doet de \omterm{def:g3:division:remainder}{rest}).

\textbf{4.} $25 \times 32 = 3\,200 \div 4 = 800$, want $25 = 100 \div 4$. En
$99 \times 47 = (100 - 1) \times 47 = 4\,700 - 47 = 4\,653$: distributiviteit
over het \omterm{ex:g1:subtraction:difference}{verschil} (\cref{thm:g7:priorities:distributivity}).

\textbf{5.} $8 \times (5 + 3) = 64$ en $8 \times 5 + 8 \times 3 = 40 + 24 = 64$:
die komen overeen --- dat is de distributiviteit. Maar
$8 \times (5 \times 3) = 8 \times 15 = 120$, terwijl
$(8 \times 5) \times (8 \times 3) = 40 \times 24 = 960$: geen overeenkomst. De
factor verdeelt zich over de \emph{termen van een \omterm{def:g1:addition:def}{som}}, niet over de factoren van
een \omterm{def:g2:mult:def}{product}.

\textbf{6.}
\begin{center}
\renewcommand{\arraystretch}{1.15}
\begin{tabular}{ccl}
$13$ & $24$ & bewaard ($13$ \omterm{def:g3:numbers:evenodd}{oneven}) \\
$6$ & $48$ & doorgestreept ($6$ \omterm{def:g3:numbers:evenodd}{even}) \\
$3$ & $96$ & bewaard \\
$1$ & $192$ & bewaard \\
\end{tabular}
\end{center}
De \omterm{def:g1:addition:def}{som} van de overgebleven rechtergetallen: $24 + 96 + 192 = 312$, en inderdaad
is $13 \times 24 = 312$.

\textbf{7.} $21$ halveren: $21, 10, 5, 2, 1$; $17$ verdubbelen:
$17, 34, 68, 136, 272$. De \omterm{def:g3:numbers:evenodd}{even} linkergetallen $10$ en $2$ worden doorgestreept;
wat overblijft is $17 + 68 + 272 = 357$. Controle: $21 \times 17 = 357$.

\textbf{8.} $16$ halveren: $16, 8, 4, 2, 1$; $25$ verdubbelen:
$25, 50, 100, 200, 400$. Elk linkergetal behalve de laatste $1$ is \omterm{def:g3:numbers:evenodd}{even}: er blijft
één rij over, en het antwoord is $400$. Een linkerkolom die uit pure
verdubbelingen bestaat, bewaart alleen haar laatste rij --- de methode valt terug
op verdubbelen alleen.

\textbf{9.} $13 = 12 + 1$, dus volgens de distributiviteit is
$13 \times 24 = 12 \times 24 + 1 \times 24$; en $12 \times 24 = 6 \times 48$,
omdat één factor halveren terwijl je de andere verdubbelt het \omterm{def:g2:mult:def}{product} gelijk laat
($12 \times 24 = 6 \times 2 \times 24 = 6 \times 48$). Het \omterm{def:g3:numbers:evenodd}{oneven} linkergetal
$13$ splitste één kopie van het rechtergetal af: de eerste rij zet $24$ op de
bank.

\textbf{10.} $6 \times 48 = 3 \times 96$ (halveren en verdubbelen, niets op de
bank: $6$ is \omterm{def:g3:numbers:evenodd}{even}). $3 \times 96 = 2 \times 96 + 96 = 1 \times 192 + 96$: de
\omterm{def:g3:numbers:evenodd}{oneven} $3$ zet $96$ op de bank. Tot slot $1 \times 192 = 192$: de \omterm{def:g3:numbers:evenodd}{oneven} $1$ zet
$192$ op de bank en de tabel is uit. Totaal op de bank:
$24 + 96 + 192 = 312$. Een rij zet haar rechtergetal precies dan op de bank als
haar linkergetal \omterm{def:g3:numbers:evenodd}{oneven} is --- het recept ``bewaar de rijen met een \omterm{def:g3:numbers:evenodd}{oneven}
linkergetal en tel op'' is dus precies de boekhouding van de distributiviteit.

\textbf{11.} $1 + 4 + 8 = 13$. Voor $21$ waren de overgebleven rijen die van
$17$, $68 = 4 \times 17$ en $272 = 16 \times 17$: de verdubbelingen
$1 + 4 + 16 = 21$.

\textbf{12.} $25 = 16 + 8 + 1$ en $42 = 32 + 8 + 2$.

\textbf{13.} Neem het grootste verdubbelgetal dat niet groter is dan het doel en
trek het af; wat overblijft is kleiner, en ook kleiner dan het net gebruikte
verdubbelgetal (anders zou een groter getal gepast hebben). Bij herhaling krimpt
de \omterm{def:g3:division:remainder}{rest} strikt en moet hij $0$ bereiken: het proces stopt, en de gebruikte
verdubbelingen zijn allemaal verschillend. Voor $100$: $64$ past ($36$ over),
$32$ past ($4$ over), $4$ past ($0$ over): $100 = 64 + 32 + 4$.

\textbf{14.} Verdubbeltabel van $35$: $1 \to 35$, $2 \to 70$, $4 \to 140$,
$8 \to 280$, $16 \to 560$. Omdat $19 = 16 + 2 + 1$, kies je die rijen:
$560 + 70 + 35 = 665$. Controle: $19 \times 35 = 665$.

\textbf{15.} Dezelfde drie als bij de schrijver en de boer:
$24 + 96 + 192$ --- de verdubbelingen van $24$ die door $13 = 1 + 4 + 8$ gekozen
worden. Getallen als sommen van verdubbelingen schrijven is precies ze met
$0$'en en $1$'en schrijven; het hele verhaal staat in
\cref{pb:g8:powers:1}.
\end{solution}
